% ===========================================================================
% CAPÍTULO 4 — RESULTADOS: caracterización eléctrica de la PCB de diseño propio
% ===========================================================================

\section{Caracterización eléctrica de la PCB de comunicación serie} \label{sec:resultados_pcb}

La comparativa de la sección anterior mide el \textbf{transporte} entre el PS y la PL, y por eso
puede ejecutarse con el lazo cerrado dentro de la FPGA. Esta sección mide otra cosa distinta: la
\textbf{placa de diseño propio}. Todo lo que se caracteriza aquí sale por los pines del MPSoC,
atraviesa los catorce transceptores RS-485/RS-422, recorre el cobre y vuelve.

La campaña se ejecuta sobre el transporte MCDMA, que es la solución final, con un
\textit{bitstream} \textbf{sin el módulo de \textit{loopback}} ---requisito que no se puede saltar:
con el \textit{bitstream} del banco de pruebas, los pines de recepción están desconectados de los
transceptores y la PCB sería sencillamente invisible.

La topología física de la placa, fijada por los \textit{jumpers}, es la siguiente:

\begin{table}[H]
    \centering
    \begin{tabular}{llll}
        \hline
        \textbf{Bus} & \textbf{Tipo} & \textbf{Maestro} & \textbf{Esclavos} \\
        \hline
        A & RS-485 multipunto (7 nodos en un bus compartido) & CH0  & CH1--CH6 \\
        B & RS-422 punto a punto (estrella desde el maestro) & CH7  & CH8--CH10 \\
        C & RS-422 punto a punto (estrella desde el maestro) & CH11 & CH12--CH13 \\
        \hline
    \end{tabular}
    \caption{Topología de buses de la PCB de comunicación serie.}
    \label{tab:pcb_topologia}
\end{table}

Es importante subrayar que \textbf{ningún test da esta topología por supuesta}: el test T2 la mide,
y el resto se apoya en lo medido. La diferencia entre el bus A y los buses B y C condiciona la
lectura de todos los resultados: en el multipunto los seis esclavos se oyen entre sí, mientras que
en los punto a punto los esclavos solo oyen al maestro.

% ---------------------------------------------------------------------------
\subsection{Los nueve tests}
% ---------------------------------------------------------------------------

\begin{table}[H]
    \centering
    \begin{tabular}{p{3.1cm} p{11.4cm}}
        \hline
        \textbf{Test} & \textbf{Resultado} \\
        \hline
        T1 Reposo & \textbf{0 bytes espurios} en los catorce canales: la polarización de \textit{fail-safe} de la placa funciona y las líneas no flotan. \\
        T2 Topología & Los tres buses están puenteados como indican los \textit{jumpers} y \textbf{no hay ni un cruce entre buses} (0 enlaces no esperados). \\
        T3 Multipunto & \textbf{6/6 esclavos} del bus A oyen íntegro al maestro, hasta 460~kbps. \\
        T4 Velocidad máx. & Bus A: \textbf{460~kbps}. Buses B y C: \textbf{1~Mbps}. Por encima, la tasa de error se dispara. \\
        T5 Limitador SLO & Con el bit SLO a 1, los tres buses llegan \textbf{limpios a 4~Mbps}. \\
        T6 Vuelta del bus & \textit{Guard time} mínimo de \textbf{0~$\mu$s}: el control de dirección del transceptor suelta la línea a tiempo y no hace falta guarda en software. \\
        T7 Diafonía & \textbf{0 bytes de fuga} del bus A hacia los buses B y C. \\
        T8 Colisión & La trama sale corrupta ---como debe--- y el bus \textbf{se recupera}: el intercambio siguiente es limpio. \\
        T9 Estabilidad & 10~s de tráfico continuo: bus A a 460~kbps, BER de 147~ppm; buses B y C a 1~Mbps, \textbf{BER de 0}. \\
        \hline
    \end{tabular}
    \caption{Resultado de los nueve tests de caracterización de la PCB.}
    \label{tab:pcb_resumen}
\end{table}

El contador de fallos de transmisión de la campaña completa es \textbf{cero}: ningún envío se quedó
sin salir, de modo que todas las cifras corresponden a bytes que de verdad viajaron por el bus. Esta
distinción no es una formalidad: confundir «el transporte no llegó a transmitir» con «el bus
corrompió los bytes» fue precisamente lo que invalidó la primera tanda de medidas, en la que los
buses B y C aparecían rotos y era mentira.

% ---------------------------------------------------------------------------
\subsection{El hallazgo principal: el limitador de slew estaba activo por defecto} \label{sec:pcb_slo}
% ---------------------------------------------------------------------------

El resultado más relevante de toda la caracterización es que \textbf{la placa venía funcionando con
el limitador de \textit{slew} de los transceptores activado sin que nadie lo supiera}, y que
desactivarlo multiplica por cuatro la velocidad de los tres buses. La medida, repetida en dos tandas
independientes, es la de la figura~\ref{fig:pcb_slo}.

\begin{figure}[H]
    \centering
    \begin{tikzpicture}
        \begin{axis}[
            benchbar,
            height=6.4cm,
            width=0.92\textwidth,
            ylabel={Velocidad máxima sin errores (Mbps)},
            ymin=0, ymax=5.2,
            bar width=22pt,
            symbolic x coords={Bus A, Bus B, Bus C},
            xtick={Bus A, Bus B, Bus C},
            nodes near coords={\pgfmathprintnumber[fixed, fixed zerofill, precision=2]{\pgfplotspointmeta}},
            legend style={at={(0.5,1.04)}, anchor=south, legend columns=2, font=\footnotesize,
                /tikz/every even column/.append style={column sep=0.5cm}},
        ]
        \addplot[fill=cGPIO,  draw=none] coordinates {(Bus A,0.46) (Bus B,1.0) (Bus C,1.0)};
        \addplot[fill=cMCDMA, draw=none] coordinates {(Bus A,4.0) (Bus B,4.0) (Bus C,4.0)};
        \legend{Limitador activo (SLO = 0\, por defecto), Limitador desactivado (SLO = 1)}
        \end{axis}
    \end{tikzpicture}
    \caption{Velocidad máxima con tasa de error nula en cada bus de la PCB, según el valor del bit
    SLO. Más es mejor. El pin \texttt{SLO} del LTC2865 es un habilitador de modo lento activo a nivel
    bajo: con el valor por defecto (0) el limitador de \textit{slew} está \textit{activo} y capa el
    driver a 250~kbps nominales, y el bus multipunto topa en 460~kbps. Desactivándolo (1), los tres
    buses llegan limpios a 4~Mbps.}
    \label{fig:pcb_slo}
\end{figure}

El resultado es contraintuitivo si uno se fía del nombre que el software da al bit: la macro
\texttt{TRANSCEIVER\_SLO\_OFF} vale 0 ---y es el valor por defecto--- y \texttt{TRANSCEIVER\_SLO\_ON}
vale 1. Poner a 1 lo que se llama «SLO activado» hace que el enlace vaya \textit{más rápido}, cuando
un limitador de \textit{slew} solo puede costar velocidad: frenar los flancos nunca puede darla.

La explicación no está en la PCB ni en el VHDL, sino en la hoja de características del transceptor.
El componente empleado es el \textbf{LTC2865} \cite{ltc2865}, y su pin \texttt{SLO} es un habilitador
de modo lento \textbf{activo a nivel bajo}:

\begin{quote}
\textit{«SLO (Slow Mode Enable): a low input switches the transmitter to the slew rate limited
250\,kbps max data rate mode. A high input supports 20\,Mbps.»}
\end{quote}

Es decir: \textbf{el hardware no está invertido, sino que se comporta exactamente como especifica su
fabricante}. Lo que está mal es el nombre de la macro en el software. Con el bit a 0 ---el valor por
defecto--- el driver queda \textbf{capado a 250~kbps nominales}, y con el bit a 1 el limitador se
\textit{desactiva} y el driver da hasta 20~Mbps. El bit no entra en el núcleo serie
---\texttt{CONFIGURABLE\_SERIAL.vhd} no tiene ninguna entrada \texttt{slo}---, sino que viaja del
registro de configuración directamente al pin del chip, sin inversores ni lógica intermedia, como
confirma el esquemático de la placa.

Las cifras medidas encajan con esa especificación con una precisión que no deja lugar a dudas:

\begin{table}[H]
    \centering
    \begin{tabular}{lrr}
        \hline
        \textbf{Bus} & \textbf{Máx. medido (SLO = 0)} & \textbf{Nominal del chip} \\
        \hline
        A (RS-485 multipunto, 7 nodos) & 460~kbps  & \multirow{3}{*}{250~kbps} \\
        B (RS-422 punto a punto)       & 1~Mbps    & \\
        C (RS-422 punto a punto)       & 1~Mbps    & \\
        \hline
    \end{tabular}
    \caption{Velocidad máxima medida con el limitador activo, frente al máximo nominal del LTC2865 en
    ese modo. Los tres buses superan el nominal, que es un valor garantizado y por tanto conservador.}
    \label{tab:pcb_slo_nominal}
\end{table}

Los tres buses aguantan \textit{por encima} de los 250~kbps nominales ---lo cual es esperable, ya que
se trata de un límite garantizado y conservador, medido aquí sobre pistas cortas y a temperatura
ambiente---, y el orden entre ellos es exactamente el que la física predice: el bus más cargado, el
multipunto de siete nodos, es el que menos margen tiene sobre el nominal, mientras que los dos
enlaces punto a punto, con la línea mucho más descargada, estiran hasta 1~Mbps. En cuanto se pide más
velocidad, el tiempo de flanco pasa a ser una fracción apreciable del tiempo de bit, el ojo se cierra
y aparecen los errores.

Conviene deshacer aquí una intuición equivocada, porque es fácil caer en ella: el limitador de
\textit{slew} \textbf{no mejora} la integridad de la señal a alta velocidad. Hace lo contrario. Su
utilidad es reducir la emisión electromagnética cuando se trabaja despacio, y ese beneficio se paga
en ancho de banda. Con el limitador desactivado el driver es de 20~Mbps, de modo que los 4~Mbps del
barrido caen \textbf{holgadamente dentro de especificación}, y de ahí que la tasa de error sea nula
en todos los puntos.

La consecuencia práctica es inmediata y notable: \textbf{la PCB da 4~Mbps en los tres buses}, cuatro
veces más de lo que se venía usando, sin tocar una sola pista de cobre y sin resintetizar nada. Basta
con poner ese bit a 1 ---y renombrar la macro, que es la que indujo el error.

La figura~\ref{fig:pcb_ber} muestra el mismo fenómeno visto desde la tasa de error del bus
multipunto, que es el más exigente de los tres.

\begin{figure}[H]
    \centering
    \begin{tikzpicture}
        \begin{axis}[
            bench,
            height=6.6cm,
            xmode=log, log basis x=10,
            xlabel={Velocidad de línea (baudios)},
            ylabel={Tasa de error, BER (ppm)},
            xmin=7000, xmax=5500000,
            ymin=-40000, ymax=1150000,
            xtick={9600, 115200, 460800, 921600, 2000000, 4000000},
            xticklabels={{9,6k}, 115k, 461k, 922k, 2M, 4M},
            scaled y ticks=false,
            ytick={0, 250000, 500000, 750000, 1000000},
            yticklabels={0, 250\,000, 500\,000, 750\,000, 1\,000\,000},
            xmajorgrids=true,
            legend style={at={(0.02,0.97)}, anchor=north west},
        ]
        \addplot[color=cGPIO, line width=1.4pt, mark=square*, mark size=2.2pt,
                 mark options={fill=cGPIO, draw=white, line width=0.6pt}]
            coordinates {(9600,0) (115200,0) (460800,0) (921600,52734) (1000000,15625) (2000000,1000000)};
        \addlegendentry{SLO = 0: limitador activo (por defecto)}

        \addplot[color=cMCDMA, line width=1.4pt, mark=*, mark size=2.4pt,
                 mark options={fill=cMCDMA, draw=white, line width=0.6pt}]
            coordinates {(9600,0) (115200,0) (460800,0) (921600,0) (1000000,0)
                         (2000000,0) (3000000,0) (4000000,0)};
        \addlegendentry{SLO = 1: limitador desactivado}

        \node[anchor=east, font=\footnotesize, text=cGPIO!85!black,
              fill=white, inner sep=1.5pt]
            at (axis cs:1900000,780000) {el enlace se cae};
        \node[anchor=west, font=\footnotesize, text=cMCDMA!80!black,
              fill=white, inner sep=1.5pt]
            at (axis cs:11000,150000) {BER 0 en todo el barrido, hasta 4~Mbps};
        \end{axis}
    \end{tikzpicture}
    \caption{Tasa de error del bus A (RS-485 multipunto, 7 nodos) frente a la velocidad de línea,
    con y sin el limitador de \textit{slew}. Con el limitador activo ---el valor por defecto--- el
    enlace se degrada a partir de 460~kbps y colapsa por completo a 2~Mbps, coherente con el máximo
    nominal de 250~kbps del chip en ese modo; desactivándolo la tasa de error es nula en todo el
    barrido. Los puntos de 3 y 4~Mbps no existen en la curva de SLO~=~0 porque el barrido se
    interrumpe en cuanto un punto no deja pasar un solo byte correcto.}
    \label{fig:pcb_ber}
\end{figure}

% ---------------------------------------------------------------------------
\subsection{Comportamiento del bus multipunto}
% ---------------------------------------------------------------------------

El bus A topa antes que los otros dos (460~kbps frente a 1~Mbps con el mismo ajuste de SLO), y esto
es coherente con su naturaleza y no un defecto de la placa: siete nodos colgados de un bus
compartido cargan más la línea y reflejan más que un enlace punto a punto. La diferencia se mantiene
en la parte alta del barrido incluso con el bit SLO a 1, de modo que es independiente del hallazgo
anterior.

El test T3 lo muestra con claridad: cuenta cuántos de los seis esclavos oyen la trama del maestro
\textbf{íntegra} a cada velocidad.

\begin{table}[H]
    \centering
    \begin{tabular}{lrrrrr}
        \hline
        \textbf{Velocidad} & 9\,600 & 115\,200 & 460\,800 & 921\,600 & 1~M \\
        \hline
        Esclavos íntegros (de 6) & 6 & 6 & \textbf{6} & 5 & 0 \\
        Esclavos parciales       & 0 & 0 & 0 & 1 & 6 \\
        \hline
    \end{tabular}
    \caption{Integridad del bus multipunto (T3) frente a la velocidad, con SLO~=~0. El escalón entre
    460~kbps y 1~Mbps es nítido.}
    \label{tab:pcb_multidrop}
\end{table}

La transición no es gradual sino un escalón: a 460~kbps los seis esclavos reciben la trama completa;
a 921~kbps ya hay uno que la recibe truncada; a 1~Mbps ninguno la recibe entera. Ese punto de
ruptura es el que define la velocidad de operación del bus, y es el que el bit SLO desplaza hasta
4~Mbps.

% ---------------------------------------------------------------------------
\subsection{Integridad, aislamiento y estabilidad}
% ---------------------------------------------------------------------------

El resto de la campaña confirma que la placa se comporta como debe en los aspectos que no dependen
de la velocidad:

\textbf{Polarización de reposo (T1).} Con los catorce transceptores callados y escuchando durante
300~ms, no aparece \textbf{ni un solo byte espurio} en ningún canal. Es la comprobación más barata
de que la red de polarización de \textit{fail-safe} existe y hace su trabajo: si la línea flotara, el
receptor amplificaría el ruido y lo entregaría como datos.

\textbf{Aislamiento entre buses (T7).} Saturando el bus A con una ráfaga a su velocidad máxima, no se
filtra \textbf{ni un byte} a los buses B y C. No hay acoplamiento apreciable entre pistas ni cruces
de cableado, lo que la matriz de conectividad de T2 confirma de forma independiente: cero enlaces no
esperados.

\textbf{Vuelta del bus (T6).} En \textit{half-duplex}, el esclavo no puede contestar hasta que el
maestro haya soltado la línea. Barriendo el tiempo de guarda desde 0 hasta 1000~$\mu$s, las ocho
respuestas llegan íntegras \textbf{ya con guarda nula}: el control de dirección que genera el propio
transceptor libera la línea a tiempo, y cualquier protocolo que se monte encima no necesita añadir
guarda por software.

\textbf{Colisión (T8).} Cuando dos esclavos emiten a la vez, la trama sale corrupta ---eso es física:
los drivers hacen un AND cableado---, pero el bus \textbf{se recupera} y el intercambio siguiente es
limpio. Ningún driver se queda enganchado conduciendo.

\textbf{Estabilidad sostenida (T9).} Diez segundos de tráfico continuo por bus, a la velocidad máxima
limpia de cada uno:

\begin{table}[H]
    \centering
    \begin{tabular}{lrrrr}
        \hline
        \textbf{Bus} & \textbf{Velocidad} & \textbf{Tramas} & \textbf{Bytes} & \textbf{BER} \\
        \hline
        A (RS-485 multipunto) & 460\,800  & 6\,781  & 433\,984 & 147 ppm \\
        B (RS-422)            & 1\,000\,000 & 14\,408 & 922\,112 & \textbf{0} \\
        C (RS-422)            & 1\,000\,000 & 14\,466 & 925\,824 & \textbf{0} \\
        \hline
    \end{tabular}
    \caption{Tasa de error en régimen permanente (T9), 10~s de tráfico continuo por bus.}
    \label{tab:pcb_estabilidad}
\end{table}

Los dos buses RS-422 no cometen \textbf{un solo error} en casi un millón de bytes cada uno. El bus
multipunto se queda en 147~ppm, una cifra baja pero no nula, coherente con la mayor exigencia de una
línea compartida por siete nodos cerca de su límite de velocidad.

% ---------------------------------------------------------------------------
\subsection{Salvedades del método}
% ---------------------------------------------------------------------------

Tres precisiones sobre lo que estos números pueden y no pueden sostener:

\textbf{Los conteos parciales de la matriz de conectividad no miden calidad de enlace.} En la matriz
de T2 aparecen enlaces con 25 de 32 bytes correctos, y sin embargo T3 y T4 a esa misma velocidad dan
seis de seis esclavos íntegros y BER nulo sobre 1024 bytes. Es un artefacto del método de T2, que
espera un tiempo fijo y lee una sola vez: si un byte llega tarde se pierde y descuadra la
comparación. T2 vale para el \textbf{mapa} de conectividad, que es para lo que está, y no para
cuantificar.

\textbf{Los ocho enlaces que faltan en la matriz no son un fallo.} De los sesenta enlaces esperados
se miden cincuenta y dos. Los ocho ausentes son los esclavos de los buses B y C, que no se oyen entre
sí porque esos buses son RS-422 punto a punto en estrella desde el maestro, y no multipunto. La
topología medida coincide exactamente con la física de la placa.

\textbf{El contador de errores de hardware de T9 no es fiable}, y no se ha usado para ninguna
conclusión: marca millones de errores en tandas cuyo BER medido es nulo, de modo que está contando
algo que no son errores de trama. Las cifras de BER de la tabla~\ref{tab:pcb_estabilidad} se calculan
comparando byte a byte lo recibido con lo enviado.
